% ============================================================
% Chapter 6: Experimental Design
%   Section: Evaluation Harness
%     Subsection: Retrieval Evaluation
% Included from chapters/06-experimental-design/04-evaluation-harness/evaluation-harness.tex
% Source: docs/decisions/evaluation.md, "Retrieval evaluation";
% retrieval.py.
% ============================================================

\subsection{Retrieval Evaluation}
\label{subsec:harness-retrieval}

RQ2 asks what retrieval contributes. The harness answers in two ways.

\paragraph{Against the annotator's evidence.} For each claim with
reference links, it checks at three depths --- every candidate seen,
the ranked evidence the model was shown, and the evidence it cited ---
whether a reference \emph{link} was found, and whether a reference
\emph{domain} was; and it computes the reciprocal rank of the first hit
in the ranked list. Links are normalised with the retriever's own
function, imported rather than copied, so ``found'' means what the
pipeline means by it. For X-Fact, links on the fact-checker's own site
are removed from the references and counted, since they are the
verdict. A miss is not proof of failure --- another source can be as
good as the annotator's --- which is why domain hits are reported beside
link hits. The cut reasons of references that were found and then cut
show where they were lost.

\paragraph{Without references.} For any set: the share of claims with
any ranked evidence, candidates and ranked sources per claim, the share
cut at ranking, distinct and rated domains, the share from the internal
corpus, retrieval and ranking time, queries per claim, and which query
kind --- anchor, proposition or refutation --- found each ranked
source. All are means over claims with bootstrap intervals, by set and
by language.

\paragraph{Strategies.} A retrieval strategy is a run with its own
thresholds, corpus, engines or code, named by a label that enters the
run key. A retrieval-only mode stops before the model, so a strategy
costs only its search and fetch time, and a comparison command sets
strategies side by side, each paired against the first. This is what
the planned ablation (Table~\ref{tab:priorities}) uses: switching a
component off on the same frozen claims, in one window, rather than
comparing versions of the code that cannot be re-run on today's web.

\paragraph{Switching queries off.} Which of the three queries are sent
is a per-run threshold: a set of query kinds, all three by default.
Only when it is not all three does it enter the run key, so every run
made before the switch keeps its key.

\paragraph{Pooled relevance judgements.} The annotator's links are a
strict witness: another page can serve as well. So, as in TREC
\cite{voorhees2005trec}, every variant's top five ranked sources per
claim are pooled, a page found by two variants entering once; a person
judges each page once, in the terminal, as bearing on the claim or not;
and every variant is scored on the same judgements --- precision at
five, recall of the pool's relevant pages, reciprocal rank, and the
share of its top five that was judged. A variant pooled after the
judging adds only its own pages and reports how much of it was not
judged.
